% ============================================================================
%  proposta-relatore.tex — documento autonomo, due pagine, da consegnare al
%  relatore per presentare il tema, l'impianto e lo stato del lavoro.
%
%  NON fa parte di main.tex: si compila da solo.
%      make proposta          (dalla radice del repository)
%  oppure:
%      cd thesis/extra && latexmk -pdf proposta-relatore.tex
%
%  Il PDF finisce in thesis/build/ e non entra in git (vedi .gitignore).
%
%  I dati identificativi arrivano da ../metadata.tex: non duplicarli qui.
%  Il titolo di questo documento e' volutamente diverso dal titolo della tesi:
%  qui serve un'esca, non l'intestazione ufficiale.
%
%  Vincoli rispettati (CLAUDE.md): nessuna data e nessuna scadenza; una frase
%  per riga; virgolette con \enquote; niente sigle interne del repository
%  (D-0NN, ADR-000N, E1..E7, V-00N) — il documento e' rivolto a chi il
%  repository non lo ha mai visto.
% ============================================================================

\documentclass[10pt,a4paper]{article}

\usepackage[T1]{fontenc}
\usepackage[utf8]{inputenc}
\usepackage[italian]{babel}
\usepackage{csquotes}
\usepackage{times}                       % togliere questa riga per tornare al Computer Modern della tesi
\usepackage[margin=2.0cm,top=1.5cm,bottom=1.4cm]{geometry}
\usepackage{microtype}
\usepackage{graphicx}
\usepackage{booktabs}
\usepackage{titlesec}
\usepackage{xcolor}
\usepackage{ragged2e}
\usepackage[hidelinks]{hyperref}

% ============================================================================
%  metadata.tex — dati identificativi della tesi.
%  Unico posto in cui modificare titolo, autore, relatore, anno accademico.
%  Ogni altro file usa questi comandi: non duplicare mai queste stringhe.
% ============================================================================

\newcommand{\tesiTitolo}{Intelligenza artificiale generativa e creativit\`a artistica: analisi di Generative Adversarial Networks e Creative Adversarial Networks}
\newcommand{\tesiTitoloBreve}{IA generativa e creativit\`a artistica}

\newcommand{\tesiAutore}{Gian}
\newcommand{\tesiMatricola}{TODO}

\newcommand{\tesiRelatore}{TODO}
\newcommand{\tesiCorrelatore}{}          % lasciare vuoto se assente

\newcommand{\tesiAteneo}{Universit\`a degli Studi di Bergamo}
\newcommand{\tesiDipartimento}{Dipartimento di Ingegneria Gestionale, dell'Informazione e della Produzione}
\newcommand{\tesiCorso}{Corso di Laurea Magistrale in Ingegneria Informatica}
\newcommand{\tesiClasse}{Classe LM-32}

\newcommand{\tesiAnnoAccademico}{2025 / 2026}

\graphicspath{{../figures/static/}}

\definecolor{acc}{HTML}{6E3030}
\definecolor{grigio}{HTML}{5F5F5F}
\definecolor{filetto}{HTML}{C8B6B6}

\titleformat{\section}{\normalfont\bfseries\color{acc}}{}{0pt}{}
\titlespacing*{\section}{0pt}{7pt}{2pt}

\setlength{\parindent}{1.2em}
\setlength{\parskip}{3pt}
\renewcommand{\arraystretch}{1.12}
\pagestyle{empty}
\hyphenpenalty=800
\emergencystretch=1.2em
\justifying

\begin{document}

% ---------------------------------------------------------------- INTESTAZIONE
\begin{center}

  {\LARGE\bfseries\color{acc} Relazione tra cretività e intelligenza artificiale generativa}\\[3pt]
  {\large Una verifica di ciò che i sistemi generativi affermano di sé}
\end{center}

\vspace{-3pt}
{\color{filetto}\rule{\linewidth}{0.8pt}}

% ---------------------------------------------------------------------- IL TEMA
\section{Il tema}

Sono partito da un'esperienza personale, non da un problema di letteratura: gli strumenti generativi permettono di produrre immagini che, per mancanza di formazione tecnica, non sono semplici da realizzare.
E questo è un beneficio reale.

Quello che non mi convince non è lo strumento: è ciò che intorno allo strumento viene affermato.
Alcune architetture non si limitano a generare immagini, ma rivendicano di essere \emph{creative}.
Le Creative Adversarial Networks sono il caso esemplare, perché la pretesa non sta in una brochure: sta nel nome e dentro la funzione di perdita.
La mia tesi vuole prendere quella rivendicazione sul serio abbastanza da sottoporla a verifica, e mostrare che cosa succede quando la si verifica davvero.

% -------------------------------------------------------------- IL DOPPIO VAGLIO
\section{Perché quella rivendicazione non regge}

La verifica procede su due piani.

Sul piano empirico, il sistema misura la propria creatività con uno strumento che si è costruito da sé: la testa di classificazione stilistica del proprio discriminatore, allenata insieme al generatore che dovrebbe giudicare.
È una parte in causa chiamata ad arbitrare la propria partita.
Ho voluto vedere che cosa succede se sugli stessi identici campioni si interroga un giudice esterno --- un classificatore addestrato una volta sola sulle sole opere reali e poi congelato.
I due responsi non si limitano a differire: si contraddicono.
Il giudice interno dichiara immagini facilmente attribuibili a uno stile; quello indipendente le trova quasi il triplo più ambigue.
La creatività dichiarata e la creatività misurata non sono la stessa grandezza, e la distanza fra le due è quantificabile.

Sul piano teorico, la rivendicazione cade prima ancora dei numeri.
Se si assume, con l'estetica simbolica di Susanne Langer, che un'opera sia arte in quanto forma imbevuta di emozione, allora ciò che manca a questi sistemi non è la potenza né la varietà: è l'intenzione emotiva che trasforma una forma in un simbolo.
Tengo a precisare che la critica non è \enquote{è solo casualità}: la letteratura sulla creatività computazionale definisce la creatività proprio \emph{in opposizione} alla casualità, e liquidare la questione così sarebbe comodo e sbagliato.
L'obiezione è più stretta, e per questo più difficile da respingere.

% ------------------------------------------------------------------- PUNTO ETICO
\section{Il punto etico, che non è un capitolo a parte}

Attribuire creatività a un sistema che non la possiede è una forma di antropomorfizzazione, e l'antropomorfizzazione non è un vezzo linguistico: distorce il giudizio morale di chi ascolta.
Ha conseguenze misurabili su come il pubblico valuta le opere, su come viene prezzato e percepito il lavoro degli artisti umani, su come si orientano mercato e regolazione.
Per questo la componente etica non si aggiunge in coda alla parte tecnica: è la stessa linea di ragionamento portata fino alle sue conseguenze.

Su un punto voglio essere netto, perché è la tensione che ho dovuto sciogliere.
Il beneficio di accessibilità da cui sono partito resta intatto, e resta \emph{indipendente} dal modo in cui il modello sottostante è stato addestrato.
Avevo pensato di difendere l'addestramento paragonandolo all'ispirazione umana --- anche un pittore impara guardando altri pittori.
Ho dovuto abbandonare l'analogia, e non per ragioni giuridiche: perfino i tribunali si sono pronunciati in modo contraddittorio su questo stesso punto.
L'ho abbandonata per coerenza con la mia stessa tesi.
Se ciò che manca alla macchina è la componente emotiva, allora manca anche nell'addestramento, non solo nell'output: e l'analogia perde il fondamento che la renderebbe una difesa.
Restano quindi due valutazioni che convivono senza annullarsi --- uno strumento può fare del bene a chi lo usa e restare, allo stesso tempo, costruito in un modo eticamente problematico.
La tesi non sceglie fra le due: le tiene insieme e distinte, perché fonderle significherebbe usare l'una per attenuare l'altra.

% ------------------------------------------------------------------ ESPERIMENTI
\section{Il lavoro sperimentale, e a che cosa serve}

L'esperimento non serve a stabilire se questi sistemi siano creativi: a quella domanda risponde il vaglio teorico.
Serve a mostrare come si comporta, sotto osservazione controllata, un sistema che rivendica una creatività che non ha --- ed è ciò che produce la misura indipendente su cui poggia tutto il resto.

Ho impostato un confronto a due condizioni con \emph{una sola} variabile indipendente: la funzione di perdita.
Generatore identico, discriminatore identico, stessi iperparametri, stessi semi casuali, stessi dati nello stesso ordine.
Le due condizioni non sono due programmi diversi: sono lo stesso programma con un interruttore.
Qualunque differenza osservata è quindi attribuibile al meccanismo di ambiguità stilistica e a nient'altro, e non a una divergenza di implementazione.
Il dataset è una collezione di trentamila opere di sei stili storici, bilanciata per costruzione; l'impianto è stato eseguito a due risoluzioni, per quattordici addestramenti complessivi.

Ogni misura è accompagnata dalla dichiarazione esplicita di ciò che \textbf{non} misura.
Gli indici standard di qualità delle immagini valutano fedeltà e varietà, non creatività, e poggiano per giunta su reti addestrate su fotografie anziché su dipinti: usarli come prova di creatività sarebbe ripetere, con altri mezzi, esattamente l'errore che la tesi contesta.

% ------------------------------------------------------------------ OSSERVAZIONI
\section{Che cosa ho osservato}

\begin{center}
  {\small
  \begin{tabular}{@{}lccl@{}}
    \toprule
     & \textbf{Baseline} & \textbf{Sistema \enquote{creativo}} & \\
    \midrule
    \multicolumn{4}{@{}l}{\textit{\color{grigio}bassa risoluzione}}\\
    Fedeltà delle immagini & 107,7 & 107,3 & indistinguibili \\
    Ambiguità stilistica   & 0,682 & 0,750 & effetto presente, gruppi separati \\
    Copertura degli stili  & 0,966 & 0,967 & nessuna deriva \\[3pt]
    \multicolumn{4}{@{}l}{\textit{\color{grigio}risoluzione raddoppiata}}\\
    Fedeltà delle immagini & 117,8 & 183,4 & degrado del 55\% \\
    Ambiguità stilistica   & 0,547 & 0,636 & effetto di entità analoga \\
    \bottomrule
  \end{tabular}}\\[3pt]
  {\footnotesize\color{grigio} Per la fedeltà il valore più basso è il migliore; per l'ambiguità il più alto.}
\end{center}

\vspace{-2pt}
Il meccanismo funziona: le immagini diventano davvero meno attribuibili a uno stile storico, e la separazione fra le due condizioni è netta.
Ma a bassa risoluzione quel guadagno non costa nulla in qualità, mentre raddoppiando la risoluzione costa più della metà della fedeltà.
È l'osservazione che considero più promettente, e la formulo con cautela perché è emersa dai dati invece di essere prevista: tutte le repliche del sistema \enquote{creativo} raggiungono il proprio massimo alla stessa epoca, molto presto, e da lì degradano insieme.

Ho trovato anche una smentita che mi è stata utile più di una conferma.
Mi aspettavo che la baseline diventasse via via più riconoscibile stilisticamente: è successo il contrario.
Una rete generativa addestrata su sei stili senza condizionamento non impara i sei stili, impara la loro miscela: più si addestra, più produce ibridi.

Due limiti li ho misurati invece di subirli, e intendo riportarli.
Il primo: il giudice indipendente non ha mai raggiunto l'affidabilità che mi ero dato come soglia, e ho verificato in tre passaggi che la causa non è la scelta degli stili ma la risoluzione --- una metrica basata su un classificatore non può essere più fine delle immagini su cui opera.
Il secondo: una replica su quattro del sistema \enquote{creativo} è degenerata, contro nessuna della baseline.
Non l'ho rimossa; è esclusa dalle medie con motivazione esplicita, perché quella differenza di stabilità è essa stessa un risultato.

\end{document}
